\documentclass[10pt,a4paper]{article}

\usepackage[margin=1.55cm]{geometry}
\usepackage[T1]{fontenc}
\usepackage[utf8]{inputenc}
\usepackage[spanish]{babel}
\usepackage{lmodern}
\usepackage{microtype}
\usepackage{graphicx}
\usepackage[dvipsnames]{xcolor}
\usepackage{tabularx}
\usepackage{array}
\usepackage{enumitem}
\usepackage{hyperref}
\usepackage{multicol}
\usepackage{tikz}
\usepackage{titlesec}
\usepackage{needspace}

\hypersetup{
  colorlinks=true,
  urlcolor=MidnightBlue,
  linkcolor=MidnightBlue,
  pdfauthor={Andrés Esteban Obando Sanipatín},
  pdftitle={CV - Andrés Esteban Obando Sanipatín}
}

\definecolor{cvblue}{HTML}{133E7C}
\definecolor{cvgray}{HTML}{50606F}
\definecolor{cvlight}{HTML}{EEF4FA}
\definecolor{cvline}{HTML}{D6E0EA}

\renewcommand{\familydefault}{\sfdefault}
\setlength{\parindent}{0pt}
\setlength{\parskip}{0.22em}
\setlist[itemize]{leftmargin=1.1em, itemsep=0.15em, topsep=0.18em}

\titleformat{\section}
  {\large\bfseries\color{cvblue}}{}{0em}{}
  [\vspace{0.2em}\color{cvline}\titlerule]
\titlespacing*{\section}{0pt}{0.7em}{0.35em}

\newcommand{\cventry}[4]{%
  \Needspace{3\baselineskip}%
  {\bfseries #1} \hfill {\small\color{cvgray} #2}\\
  {\itshape #3}\\[-0.25em]
  #4\par\vspace{0.25em}
}

\newcommand{\cvtag}[1]{%
  \tikz[baseline=(txt.base)]\node[rounded corners=4pt, fill=cvlight, text=cvblue, inner xsep=5pt, inner ysep=3pt] (txt) {\small\strut #1};%
}

\newcommand{\roundphoto}[2]{%
\begin{tikzpicture}
  \clip (0,0) circle (#1);
  \node at (0,0) {\includegraphics[width=#2]{Yo_CV_crop.jpg}};
  \draw[cvblue, line width=1pt] (0,0) circle (#1);
\end{tikzpicture}%
}

\begin{document}
\pagestyle{empty}

\begin{tabularx}{\textwidth}{@{}m{2.8cm}@{\hspace{0.55cm}}X@{}}
\centering\roundphoto{1.25cm}{3.0cm}
& {\fontsize{19}{21}\selectfont\bfseries\color{cvblue} ANDRÉS ESTEBAN OBANDO SANIPATÍN}\\[0.15em]
& {\normalsize\color{cvgray} Ingeniero Químico \textbar{} Docente Universitario \textbar{} Máster en Energías Renovables}\\[0.3em]
& {\small\color{cvgray} Quito, Ecuador \textbar{} 0984498811 \textbar{} \href{mailto:aeosangabriel@hotmail.com}{aeosangabriel@hotmail.com}}\\
& {\small\color{cvgray} \href{mailto:aeobandos@uce.edu.ec}{aeobandos@uce.edu.ec} \textbar{} \href{https://www.linkedin.com/in/andresobando1101}{linkedin.com/in/andresobando1101}}\\
\end{tabularx}

\section*{Perfil profesional}
Ingeniero químico con formación de cuarto nivel en \textbf{Energías Renovables} y \textbf{Docencia Superior Universitaria}. Experiencia en docencia universitaria, nivelación académica, apoyo a laboratorios, investigación aplicada y análisis de datos. Actualmente se desempeña como \textbf{Docente Auxiliar 1} de la Universidad Central del Ecuador, donde imparte \textbf{Química General I} y \textbf{Balance de Energía} en la carrera de Ingeniería Química. Su perfil integra base científica, enfoque pedagógico, manejo de herramientas digitales y experiencia en procesos académicos e institucionales.

\section*{Experiencia profesional}
\cventry{Docente Auxiliar 1}{Abr. 2026 -- Actualidad}{Universidad Central del Ecuador \textbar{} Facultad de Ingeniería Química}{%
\begin{itemize}
  \item Docente de las asignaturas \textbf{Química General I} y \textbf{Balance de Energía} en la carrera de Ingeniería Química.
  \item Planificación, acompañamiento académico y evaluación del aprendizaje en contexto universitario.
\end{itemize}}

\cventry{Técnico Docente}{Jul. 2024 -- Abr. 2026}{Universidad Central del Ecuador \textbar{} Facultad de Ingeniería Química}{%
\begin{itemize}
  \item Impartición de cursos de nivelación y apoyo académico en Física, Matemática y Geometría.
  \item Soporte a laboratorios, actividades académicas, investigación y procesos administrativos de la Facultad.
\end{itemize}}

\cventry{Analista de Datos}{Nov. 2023 -- Mar. 2024}{ESPE Innovativa}{%
\begin{itemize}
  \item Participación en el proyecto \textbf{RENAGRO}, incluyendo estructuración de datos y apoyo en formularios digitales para captura y tratamiento de información.
\end{itemize}}

\cventry{Docente de nivelación y apoyo académico}{Feb. 2022 -- Jun. 2024}{Grupo Impulsar \textbar{} Preuniversitario Politécnica}{%
\begin{itemize}
  \item Nivelación y clases en Física, Química, Matemática, Geometría y razonamiento lógico.
\end{itemize}}

\cventry{Docente}{Nov. 2020 -- Abr. 2022}{Instituto Superior Tecnológico Superarse}{%
\begin{itemize}
  \item Docencia en Inglés Básico, Inglés A1 e Higiene de los Alimentos.
\end{itemize}}

\cventry{Experiencia en investigación y laboratorio}{2020 -- 2023}{Escuela Politécnica Nacional \textbar{} Agroindustria LCM}{%
\begin{itemize}
  \item Investigación en extracción de oleorresinas, ensayos electroquímicos, espectroscopía Raman y caracterización experimental.
  \item Verificación de procesos, BPM y control de calidad en producción alimentaria.
\end{itemize}}

\newpage

\section*{Formación académica}
\cventry{Máster Universitario en Energías Renovables}{Dic. 2025}{Universidad Internacional de La Rioja (UNIR)}{}
\cventry{Máster Universitario en Docencia Superior Universitaria}{Abr. 2024}{Universidad Internacional de La Rioja (UNIR)}{}
\cventry{Ingeniero Químico}{Abr. 2020}{Escuela Politécnica Nacional}{}
\cventry{Productor Musical}{May. 2023}{Instituto Superior Tecnológico de Artes Visuales}{}

\section*{Certificaciones seleccionadas}
\begin{multicols}{2}
\begin{itemize}
  \item Programación en Python y Manejo de Datos (UCE, 30 h, 2025)
  \item Introducción a la normativa ISO/IEC 17025 (UCE-FIQ, 10 h, 2025)
  \item Aplicación de la Inteligencia Artificial en la Educación Superior (UCE, 20 h, 2024)
  \item Metodología de la Investigación (UCE, 20 h, 2024)
  \item Fundamentos de Contratación Pública (SERCOP, 2024)
  \item Formación de Formadores (certificación por competencias laborales)
\end{itemize}
\end{multicols}

\section*{Competencias y áreas de trabajo}
\vspace{0.15em}
\cvtag{Docencia universitaria}
\hspace{0.35em}
\cvtag{Química General I}
\hspace{0.35em}
\cvtag{Balance de Energía}
\hspace{0.35em}
\cvtag{Investigación aplicada}\\[0.35em]
\cvtag{Apoyo a laboratorios}
\hspace{0.35em}
\cvtag{Análisis de datos}
\hspace{0.35em}
\cvtag{Energías renovables}
\hspace{0.35em}
\cvtag{Innovación educativa}\\[0.35em]
\cvtag{Herramientas de IA}
\hspace{0.35em}
\cvtag{Control de calidad}
\hspace{0.35em}
\cvtag{Procesos académicos}

\vspace{0.55em}
\begin{tabularx}{\textwidth}{@{}>{\bfseries\color{cvblue}}p{4.2cm}X@{}}
Herramientas digitales & Office 365, Google Workspace, Moodle, Teams, SPSS, CSPro, Statgraphics, Python, ChatGPT, Gemini y Claude. \\
Software creativo & Adobe Illustrator, Adobe Premiere, Photoshop, Pro Tools, Ableton Live y Reaper. \\
Fortalezas profesionales & Planificación docente, comunicación académica, trabajo en equipo, organización, adaptación tecnológica, redacción técnica y soporte institucional. \\
Líneas de interés & Docencia en ingeniería química, energías renovables, corrosión y bioinhibidores, análisis de datos, metodologías de investigación y uso pedagógico de inteligencia artificial. \\
Idiomas & Español nativo; Inglés nivel B1. \\
\end{tabularx}

\section*{Proyectos y actividad complementaria}
\begin{itemize}
  \item Participación en proyectos vinculados con \textbf{inhibidores de corrosión}, caracterización de materiales y estudio de \textbf{oleorresinas}.
  \item Experiencia complementaria en \textbf{producción musical} para videojuegos, composición sonora y trabajo con herramientas digitales creativas.
  \item Perfil con capacidad para articular docencia, investigación, análisis de datos y recursos tecnológicos en entornos educativos y académicos.
\end{itemize}

\end{document}
